% code-contact-map-generator.tex
% Code Contact Map Generator: A Measure-Theoretic Framework for Minimum
% Sufficient Mutual Individuation in Software Systems
%
% Compile with: pdflatex -> bibtex -> pdflatex -> pdflatex

\documentclass[12pt,a4paper,twocolumn]{article}

% --------------------------------------------------------------------------
% Packages
% --------------------------------------------------------------------------
\usepackage[a4paper, margin=2cm, columnsep=0.8cm]{geometry}
\usepackage{amsmath,amssymb,amsthm}
\usepackage{mathtools}
\usepackage{bbm}
\usepackage{stmaryrd}        % \llbracket, \rrbracket
\usepackage{graphicx}
\usepackage{xcolor}
\usepackage{hyperref}
\usepackage{natbib}
\usepackage{listings}
\usepackage{microtype}
\usepackage{booktabs}
\usepackage{enumitem}
\usepackage{float}
\usepackage{caption}
\usepackage{subcaption}

% --------------------------------------------------------------------------
% Hyperref setup
% --------------------------------------------------------------------------
\hypersetup{
  colorlinks=true,
  linkcolor=blue!60!black,
  citecolor=green!50!black,
  urlcolor=blue!70!black,
  pdftitle={Code Contact Map Generator},
  pdfauthor={},
}

% --------------------------------------------------------------------------
% Listings / DSL syntax
% --------------------------------------------------------------------------
\definecolor{dsl-keyword}{RGB}{0,100,200}
\definecolor{dsl-string}{RGB}{160,40,40}
\definecolor{dsl-comment}{RGB}{80,120,80}
\definecolor{dsl-number}{RGB}{140,60,0}
\definecolor{dsl-bg}{RGB}{248,248,252}
\definecolor{dsl-frame}{RGB}{200,200,220}

\lstdefinelanguage{WindTunnel}{
  morekeywords={scope,analyse,compose,assert,recurse,report,
                target,exclude,method,model,query,merge,resolve,
                extract,claim,on_fail,output,format,until,while,
                static,dynamic,semantic,hf,individuation,regime,
                distance,narrow_scope,marginal_return,floor,by},
  sensitive=true,
  morestring=[b]",
  morestring=[b]',
  morecomment=[l]{\#},
}

\lstset{
  language=WindTunnel,
  backgroundcolor=\color{dsl-bg},
  frame=single,
  framerule=0.4pt,
  rulecolor=\color{dsl-frame},
  basicstyle=\ttfamily\footnotesize,
  keywordstyle=\bfseries\color{dsl-keyword},
  stringstyle=\color{dsl-string},
  commentstyle=\itshape\color{dsl-comment},
  numberstyle=\tiny\color{gray},
  numbers=left,
  numbersep=5pt,
  breaklines=true,
  breakatwhitespace=false,
  captionpos=b,
  tabsize=2,
  showspaces=false,
  showstringspaces=false,
}

% --------------------------------------------------------------------------
% Theorem environments
% --------------------------------------------------------------------------
\theoremstyle{plain}
\newtheorem{theorem}{Theorem}[section]
\newtheorem{lemma}[theorem]{Lemma}
\newtheorem{proposition}[theorem]{Proposition}
\newtheorem{corollary}[theorem]{Corollary}

\theoremstyle{definition}
\newtheorem{definition}[theorem]{Definition}
\newtheorem{example}[theorem]{Example}

\theoremstyle{remark}
\newtheorem{remark}[theorem]{Remark}

% --------------------------------------------------------------------------
% Notation macros
% --------------------------------------------------------------------------
\newcommand{\Borel}{\mathcal{B}}
\newcommand{\Part}{\mathcal{P}}
\newcommand{\Sep}{\mathrm{Sep}}
\newcommand{\CM}{\mathrm{CM}}
\newcommand{\bmin}{\beta_{\min}}
\newcommand{\bP}{\beta_{\mathcal{P}}}
\newcommand{\Scoord}{S}
\newcommand{\dS}{d_S}
\newcommand{\Neg}{\mathrm{Neg}}
\newcommand{\Rest}{\mathrm{R}_{\mathrm{est}}}
\newcommand{\Lmin}{L_{\min}}
\newcommand{\Smin}{S_{\min}}
\newcommand{\CMbmin}{\CM_{\bmin}}
\newcommand{\GC}{G_C}
\newcommand{\GS}{G_S}
\newcommand{\hol}{\mathrm{hol}}
\newcommand{\diag}{\mathrm{diag}}
\newcommand{\val}{\mathrm{val}}
\newcommand{\Sbmin}{S_{\bmin}}
\newcommand{\Real}{\mathbb{R}}
\newcommand{\Nat}{\mathbb{N}}
\newcommand{\BRS}{(\Omega,d,\mu)}

% --------------------------------------------------------------------------
% Title
% --------------------------------------------------------------------------
\title{\textbf{Code Contact Map Generator}\\[4pt]
  {\large A Measure-Theoretic Framework for Minimum Sufficient\\
  Mutual Individuation in Software Systems}}

\author{Wind Tunnel Research\\[4pt]
  \small \texttt{wind-tunnel.dev}}

\date{June 2026}

% ==========================================================================
\begin{document}
% ==========================================================================

\maketitle

% --------------------------------------------------------------------------
\begin{abstract}
% --------------------------------------------------------------------------
We introduce the \emph{Code Contact Map Generator}, a framework for
characterising software systems not by the binary outcomes of individual test
assertions, but by the \emph{minimum sufficient contact map} --- the minimal
measure-theoretic structure that certifies mutual individuation across every
pair of system components.  The foundational object is the \emph{bounded
resolvable space} $\BRS$, a compact metric space equipped with a finite
non-atomic Borel measure and a resolution floor.  Within this space every
valid partition of the system into components must deposit a non-zero
``separator'' between each pair of adjacent components: the \emph{partition
floor} $\bmin > 0$ is strictly positive and cannot be driven to zero by any
finite investigation strategy.  We prove three independent derivations of this
floor (geometric, representational, economic) and derive from it the
\emph{Minimum Sufficiency Theorem}: the partition floor $\bmin$, the minimum
sufficient contact map $\CMbmin$, the minimum sufficient Lagrangian $\Lmin$,
and the minimum sufficient action $\Smin$ are in canonical bijective
correspondence.  Every additional structural annotation (type declarations,
test suites, runtime checks) is a \emph{partition depth inflation} of
$\CMbmin$; none reduces the floor to zero.  We then model software systems
as bounded resolvable spaces, define S-entropy coordinates for code units, and
use them to derive five coordination regimes quantified by the static order
parameter $\Rest$.  The framework is realised by the \emph{Wind Tunnel DSL},
a partition-specification language in which both users and the contact-map
generator co-author investigation scripts.  The DSL is specified by a formal
context-free grammar with compositional, monotone, floor-aware operational
semantics.  We conclude that the classical test paradigm --- asking
``pass/fail?'' at a point --- is strictly weaker than contact map generation,
which asks for the minimum sufficient structure of mutual individuation across
the whole system.
\end{abstract}

\tableofcontents

\onecolumn   % switch to single column for the body (easier for long math)

% ==========================================================================
\section{Introduction}
% ==========================================================================

Software verification has, since its inception, been organised around a
single epistemic act: the evaluation of a predicate at a point.  A test is
executed, a return value is inspected, a property is declared true or false.
This paradigm is so deeply embedded in engineering practice that it shapes not
only tooling but conceptual vocabulary: ``passing the tests,'' ``green build,''
``coverage report.''  The implicit model is that correctness is a point in a
binary space, and that enough point evaluations constitute a proof.

This paper challenges that model at its foundation.  We argue that the
natural unit of software-system characterisation is not the truth value of an
isolated assertion but the \emph{contact map}: the assignment, to each pair of
system components, of the minimum information required to distinguish them from
each other.  A test failure is evidence that contributes to the contact map;
it is not the map itself.  The contact map is a global structure; point
evaluations are local observations.  The relationship is analogous to the
distinction between a measure and the value of its density at a single point.

\subsection{The Problem with Truth-at-a-Point}

Let $f : X \to Y$ be a software function and $\phi : X \times Y \to
\{0,1\}$ a specification predicate.  Classical testing evaluates $\phi(x_i,
f(x_i))$ for a finite sample $\{x_i\}_{i=1}^N \subset X$.  Several
fundamental obstacles limit the conclusions available from this procedure.

\begin{enumerate}[label=(\roman*)]
\item \textbf{Rice's theorem} \citep{rice1953classes}: no non-trivial semantic
  property of a program is decidable.  Any complete automated oracle is
  impossible.

\item \textbf{Undecidability of halting} \citep{turing1936computable}: the
  test harness itself cannot determine, in general, whether the system under
  test terminates.

\item \textbf{Local/global gap}: $\phi(x_i, f(x_i)) = 1$ for all sampled
  $x_i$ does not imply $\forall x \in X\!: \phi(x,f(x)) = 1$.  This is not
  merely a practical limitation; it is a categorical one.

\item \textbf{Compositionality failure}: local correctness of each component
  does not imply global correctness of the composed system (formalised in
  Theorem~\ref{thm:local-invisibility} below).
\end{enumerate}

These obstacles are well known individually.  What is less well recognised is
that they share a common source: the point-evaluation paradigm attempts to
characterise a measure-theoretic object (the behaviour of $f$ over a
continuous or large discrete domain) by a finite sequence of point samples.
The information deposited by each sample is bounded; the object is not.

\subsection{Contact Map Generation as an Alternative}

We propose replacing the question ``did this unit pass or fail?'' with the
question ``what is the minimum sufficient structure of mutual individuation
across this system?''  These are categorically different questions.  The
first is local and binary; the second is global and metric.  The framework
developed in this paper provides:

\begin{enumerate}[label=(\roman*)]
\item A \textbf{mathematical foundation} (Sections~\ref{sec:foundations}
  and~\ref{sec:mscm}) based on measure theory and compact metric spaces,
  yielding the \emph{partition floor} $\bmin > 0$ as the irreducible
  informational residue of any finite investigation.

\item A \textbf{model of software systems} (Section~\ref{sec:software}) as
  bounded resolvable spaces, with S-entropy coordinates for code units and
  five quantified coordination regimes.

\item A \textbf{contact map generator} (Section~\ref{sec:generator}) that
  uses static analysis, dynamic tracing, and semantic inference as partition
  operations contributing estimates of $\CMbmin$.

\item The \textbf{Wind Tunnel DSL} (Section~\ref{sec:dsl}), a partition-
  specification language with formal grammar and compositional operational
  semantics, in which both users and the system co-author investigation
  scripts.
\end{enumerate}

\subsection{Overview and Contributions}

The central result is the \emph{Minimum Sufficiency Theorem}
(Theorem~\ref{thm:minimum-sufficiency}), which establishes a canonical
bijective correspondence between the partition floor, the minimum sufficient
contact map, the minimum sufficient Lagrangian, and the minimum sufficient
action.  All additional structural annotations are \emph{partition depth
inflations} of this ground state.

The Wind Tunnel DSL (Section~\ref{sec:dsl}) is the language in which both
questions --- classical test assertions and contact map claims --- can be
expressed.  Only the second is the native question of the framework.  A DSL
execution file is a complete, reproducible record of a contact map
investigation, suitable for auditing, extension, and collaborative authorship
between human and automated agents.

% ==========================================================================
\section{Mathematical Foundations}
\label{sec:foundations}
% ==========================================================================

\subsection{Bounded Resolvable Spaces}

\begin{definition}[Bounded Resolvable Space]
\label{def:brs}
A \emph{bounded resolvable space} is a triple $\BRS$ where:
\begin{enumerate}[label=(\roman*)]
  \item $(\Omega, d)$ is a compact metric space;
  \item $\mu$ is a finite non-atomic Borel measure on $\Omega$ with
    $\mu(\Omega) < \infty$;
  \item There exist constants $\delta, c_0 > 0$ and $n \geq 1$ such that
    for every $x \in \Omega$ and every $0 < r < \delta$,
    \begin{equation}
      \mu\!\left(B(x,r)\right) \;\geq\; c_0\, r^n,
      \label{eq:resolution-floor}
    \end{equation}
    where $B(x,r) = \{y \in \Omega : d(x,y) < r\}$ is the open $d$-ball.
\end{enumerate}
The triple $(\delta, c_0, n)$ is called the \emph{resolution data}, and
condition~\eqref{eq:resolution-floor} is the \emph{resolution floor}.
\end{definition}

\begin{remark}
The resolution floor is a quantitative lower Ahlfors regularity condition
\citep{mattila1995geometry, falconer2003fractal}: every ball of radius $r$
carries at least $c_0 r^n$ measure.  It excludes purely singular or purely
atomic measures and ensures that every open set of positive diameter carries
positive measure.  Euclidean space with Lebesgue measure satisfies this with
$n = \dim$; self-similar fractals satisfying the open set condition satisfy
it with $n$ equal to the Hausdorff dimension.
\end{remark}

\subsection{Partitions and Partition Depth}

\begin{definition}[Valid Partition]
\label{def:partition}
A \emph{valid partition} of $\BRS$ is a finite collection
$\Part = \{A_1, \ldots, A_k\}$ of measurable subsets of $\Omega$ satisfying:
\begin{enumerate}[label=(\roman*)]
  \item \textbf{Covering}: $\mu\!\left(\Omega \setminus \bigcup_{i=1}^k A_i\right) = 0$;
  \item \textbf{Essential disjointness}: $\mu(A_i \cap A_j) = 0$ for $i \neq j$;
  \item \textbf{Connectedness}: each $A_i$ is connected in the subspace
    topology induced by $d$;
  \item \textbf{Positive measure}: $\mu(A_i) > 0$ for all $i$.
\end{enumerate}
Denote the set of all valid partitions by $\mathfrak{P}(\Omega,d,\mu)$.
\end{definition}

\begin{definition}[Separator]
\label{def:separator}
For $A \in \Part$ with $\Part \in \mathfrak{P}$, the \emph{separator set}
of $A$ in $\Part$ is
\begin{equation}
  \Sep(A, \Part) \;=\; \bigl\{ B \in \Part \setminus \{A\}
    : \overline{A} \cap \overline{B} \neq \emptyset \bigr\},
\end{equation}
the collection of components whose topological closures share a boundary
point with $\overline{A}$.
\end{definition}

\begin{definition}[Partition Depth]
\label{def:partition-depth}
The \emph{partition depth} of $A \in \Part$ is
\begin{equation}
  \bP(A) \;=\; \inf_{B \in \Sep(A,\Part)} \mu(B).
\end{equation}
The \emph{partition floor} of the space is
\begin{equation}
  \bmin \;=\; \inf_{\Part \in \mathfrak{P}}\; \inf_{A \in \Part}\; \bP(A).
\end{equation}
\end{definition}

\begin{remark}
Intuitively, $\bP(A)$ is the smallest measure carried by any component
adjacent to $A$.  The partition floor $\bmin$ is the infimum over all possible
partitions of this quantity: the smallest separator any partition strategy can
produce.
\end{remark}

\subsection{Positivity of the Partition Floor}

\begin{theorem}[Partition Floor is Strictly Positive]
\label{thm:floor-positive}
Let $\BRS$ be a bounded resolvable space.  Then $\bmin > 0$.
\end{theorem}

We give three independent proofs.

\paragraph{Proof I: Geometric.}

Let $\Part \in \mathfrak{P}$ be any valid partition and $B \in \Part$ any
component.  Since $B$ is connected and has positive measure, $B$ contains an
open set $U \subseteq B$ with $\mu(U) > 0$ \citep{rudin1987real}.  By
compactness of $\overline{B}$ and the resolution floor~\eqref{eq:resolution-floor},
for any $x \in U$ there exists $r_x > 0$ such that $B(x, r_x) \subseteq U$
and
\begin{equation}
  \mu(B) \;\geq\; \mu(U) \;\geq\; \mu\!\left(B(x,r_x)\right)
    \;\geq\; c_0\, r_x^n \;>\; 0.
\end{equation}
Since $k = |\Part| < \infty$ and each component has positive measure,
\begin{equation}
  \inf_{B \in \Part} \mu(B) \;\geq\;
    \frac{\mu(\Omega)}{k} \cdot \frac{1}{k} \;>\; 0,
\end{equation}
and since this holds for every $\Part \in \mathfrak{P}$, and any sequence of
partitions with $\bP(A_i) \to 0$ would require $|\Part_i| \to \infty$ in a
manner that contradicts the compactness of $\Omega$ and the atomless
condition, we conclude $\bmin > 0$.  $\square$

\paragraph{Proof II: Representational.}

Suppose $\Omega$ is uncountably infinite (which holds in all cases of
interest: $\Omega \cong [0,1]^m$ or a fractal subset thereof).  A valid
partition $\Part$ with $|\Part| = k$ cells can name at most $2^k$ distinct
regions via Boolean combinations of cells.  There are uncountably many Borel
subsets of $\Omega$ \citep{bogachev2007measure}; hence any finite partition
leaves uncountably many distinctions unrepresented.  If $\bP(A) \to 0$, then
arbitrarily fine partitions would be required, but each partition is finite
by Definition~\ref{def:partition}.  The representational capacity of a
$k$-cell partition is bounded by $k \log k$ bits of structural information
\citep{cover2006elements}, while the information content of $\Omega$ is
infinite.  Hence for any fixed finite $k$, the partition floor is bounded
below by the per-cell measure floor $c_0 (\mu(\Omega)/k)^{n/m}$, which is
positive.  $\square$

\paragraph{Proof III: Economic.}

Let $g : (0,\infty) \to (0,\infty)$ be a cost function satisfying $g(t) \to
+\infty$ as $t \to 0^+$.  Each partition operation of depth $t$ incurs cost
$g(t)$.  The return of any partition operation is bounded by $\mu(\Omega) <
\infty$ (the total available measure to be partitioned).  For an agent with
finite total budget $C < \infty$, the sequence of partition operations
$(t_1, t_2, \ldots)$ satisfies $\sum_j g(t_j) \leq C$.  Since $g(t) \to
+\infty$ as $t \to 0^+$, only finitely many $t_j$ can approach $0$.  The
agent therefore halts at some $t^* > 0$, which serves as a lower bound on
the achievable partition depth.  This bound is uniform over all finite
strategies, yielding $\bmin > 0$.  $\square$

\subsection{Individuation by Negation}

\begin{definition}[Negation of a Component]
\label{def:negation}
For $A \in \Part$, the \emph{negation} of $A$ in $\Part$ is
\begin{equation}
  \Neg(A, \Part) \;=\; \Omega \setminus A \;=\; \bigcup_{B \in \Part, B \neq A} B.
\end{equation}
\end{definition}

\begin{proposition}[Selector-Free Individuation]
\label{prop:selector-free}
The negation $\Neg(A,\Part)$ is the unique individuating complement of $A$
that requires no distinguished element or selector function.  Specifically,
$A$ is the $\mu$-a.e.\ unique set satisfying $A \cup \Neg(A,\Part) = \Omega$
and $\mu(A \cap \Neg(A,\Part)) = 0$.
\end{proposition}

\begin{proof}
By essential disjointness (Definition~\ref{def:partition}(ii)) and the
covering condition (Definition~\ref{def:partition}(i)), the complement
$\Omega \setminus A$ is well-defined up to a $\mu$-null set.  No choice
function is needed: the complement is simply the union of all other cells.
Uniqueness follows from the atomless property of $\mu$: if $A'$ also satisfies
$A' \cup (\Omega \setminus A) = \Omega$ and $\mu(A' \cap (\Omega \setminus
A)) = 0$, then $\mu(A \triangle A') = 0$, so $A = A'$ $\mu$-a.e.
\end{proof}

\begin{remark}
This proposition formalises the philosophical principle that individuation
proceeds by negation: a component $A$ is fixed as ``the whole minus everything
it is not.''  The complement carries no distinguished element; it is the
structural residue of $A$'s existence within $\Part$.
\end{remark}

\subsection{Detectability}

\begin{definition}[Signal Measure]
\label{def:signal}
For a component $A \in \Part$, its \emph{signal measure} is
\begin{equation}
  \sigma(A) \;=\; \mu(A) - \mu(\partial A),
\end{equation}
where $\partial A = \overline{A} \setminus \mathrm{int}(A)$ is the topological
boundary of $A$ in $(\Omega, d)$.
\end{definition}

\begin{definition}[Detectability]
\label{def:detectability}
A component $A \in \Part$ is \emph{detectable} if
\begin{equation}
  \sigma(A) \;>\; \bP(A).
\end{equation}
It is \emph{subdetectable} (engulfed by its own boundary) if
$\sigma(A) \leq \bP(A)$.
\end{definition}

\begin{proposition}
In a bounded resolvable space, a component $A$ is detectable if and only if
it can be distinguished, with positive $\mu$-probability, from any element
of $\Sep(A,\Part)$ by a measure-preserving observation process.
\end{proposition}

\begin{proof}
If $\sigma(A) > \bP(A)$, then $A$ carries strictly more interior measure
than the measure of any adjacent separator.  A measure-preserving process
visiting $A$ will encounter its interior --- not merely its boundary --- with
positive probability, enabling distinction.  Conversely, if $\sigma(A) \leq
\bP(A)$, then the boundary accounts for at least as much measure as any
adjacent component; observations of $A$ are indistinguishable from
observations of $\partial A$, which is shared with adjacent components.
\end{proof}

\subsection{Locating vs.\ Visiting Asymmetry}

\begin{definition}[Measure-Preserving Traversal]
\label{def:traversal}
A \emph{measure-preserving traversal} of $\BRS$ is a measure-preserving
ergodic transformation $T : \Omega \to \Omega$.  For a component $A$,
the \emph{visiting frequency} of $A$ under $T$ is
\begin{equation}
  \nu_T(A) \;=\; \lim_{N\to\infty} \frac{1}{N} \sum_{j=0}^{N-1}
    \mathbf{1}_A\!\left(T^j(x)\right) \;=\; \mu(A),
\end{equation}
where the last equality holds $\mu$-a.e.\ by the Birkhoff ergodic theorem
\citep{halmos1950measure}.
\end{definition}

\begin{definition}[Locating Cost]
\label{def:locating-cost}
The \emph{locating cost} of a component $A$ with visiting probability
$p_A = \mu(A)/\mu(\Omega)$ is
\begin{equation}
  \lambda(A) \;=\; \bmin \cdot \left\lceil \log_2 \frac{1}{p_A} \right\rceil.
\end{equation}
\end{definition}

\begin{theorem}[Locating/Visiting Asymmetry]
\label{thm:locating-visiting}
Let $T$ be a measure-preserving traversal of $\BRS$ and $A \in \Part$ a
component with $p_A = \mu(A)/\mu(\Omega)$.  Then:
\begin{enumerate}[label=(\roman*)]
  \item $T$ visits $A$ with frequency $p_A$ regardless of how small $p_A$
    is, and the visiting frequency entails no additional cost beyond the
    floor $\bmin$ already deposited by the partition.
  \item Locating $A$ (certifying that a given observation is in $A$ rather
    than in some $B \in \Sep(A,\Part)$) costs at least $\lambda(A) > 0$.
  \item As $p_A \to 0$, visiting frequency $\to 0$ but locating cost
    $\to +\infty$.  These are categorically distinct relations.
\end{enumerate}
\end{theorem}

\begin{proof}
(i) Follows directly from the Birkhoff ergodic theorem applied to
$\mathbf{1}_A$ \citep{halmos1950measure}.  The ergodic average converges to
$\mu(A) = p_A \cdot \mu(\Omega)$.  No additional information beyond
$\Part$ is required to perform the traversal; the partition was already
deposited at cost $\geq \bmin$.

(ii) To locate $A$, one must distinguish it from all $B \in \Sep(A,\Part)$,
of which there are at most $|\Part| - 1$.  Each binary decision requires at
least $\bmin$ information (by the partition floor theorem).  A binary search
tree over the $\lceil\log_2(1/p_A)\rceil$ bits needed to identify $A$ within
the partition structure therefore costs at least $\lambda(A) = \bmin \cdot
\lceil\log_2(1/p_A)\rceil$.

(iii) As $p_A \to 0$, the Birkhoff average converges to $0$ (the set is
visited rarely), while $\lceil\log_2(1/p_A)\rceil \to +\infty$ (exponentially
many partition decisions are needed to isolate it).  The visiting relation is
linear in $p_A$; the locating cost is logarithmic in $1/p_A$ times the floor
$\bmin$.  These are formally distinct: no re-parameterisation of one yields the
other.
\end{proof}

\subsection{Non-Self-Verifiability}

\begin{theorem}[Non-Self-Verifiability]
\label{thm:non-self-verify}
There is no total computable function $V : \mathfrak{P} \to \{0,1\}$ such
that $V(\Part) = 1$ if and only if $\Part$ is a valid partition and $\bP(A)
\geq \bmin$ for all $A \in \Part$.
\end{theorem}

\begin{proof}
We use a diagonal argument in the spirit of Lawvere's fixed-point theorem
\citep{lawvere1969diagonal}.  Suppose for contradiction that such a $V$
exists.  Consider the family of partitions $\{\Part_i\}$ indexed by programs
$p_i$ in a universal Turing machine.  Define a partition $\Part^*$ via a
diagonalisation:
\begin{equation}
  \Part^* = \mathrm{diag}\!\left(\{\Part_i\}\right)
    \;=\; \left\{\, A_i^* \;:\; \bP(A_i^*) \neq \bP(A_i)\right\}.
\end{equation}
If $V$ were total, then $V(\Part^*)$ would be computable.  But $V(\Part^*)
= 1$ implies $\Part^* \in \{\Part_i\}$ for some $i$, and the definition of
$\Part^*$ then implies $\bP(A_i^*) \neq \bP(A_i)$, a contradiction.  Hence
$V(\Part^*) = 0$, which means $V$ fails to verify a valid partition.  The
diagonal argument shows no total $V$ can be both complete and consistent.

More concretely: to verify $\bP(A) \geq \bmin$ for $A$, one must decide $A$
against $\Neg(A,\Part)$ --- a self-referential computation.  By
G\"odel's incompleteness theorem \citep{godel1931formal}, any sufficiently
expressive formal system contains true statements that cannot be proved within
the system.  The irreducible residue $\bmin$ corresponds to such a statement:
it is the floor that cannot be exhibited as a locatable set within the
verifier's own language.
\end{proof}

\subsection{Non-Return Theorem}

\begin{definition}[Committed Record]
\label{def:committed-record}
Let $(\Part^{(t)})_{t \geq 0}$ be a sequence of partitions produced by a
contact map investigation.  The \emph{committed record} at time $t$ is
\begin{equation}
  M(t) \;=\; \sum_{s=0}^{t} \sum_{A \in \Part^{(s)}} \mathbf{1}[\bP^{(s)}(A)
    \text{ is committed}],
\end{equation}
the total count of partition depths that have been committed (written to a
permanent record).
\end{definition}

\begin{theorem}[Non-Return]
\label{thm:non-return}
The committed record $M(t)$ is monotone non-decreasing: $M(t+1) \geq M(t)$
for all $t \geq 0$.  No investigation process can return to a prior
individuation state once a partition depth has been committed.
\end{theorem}

\begin{proof}
Commitment is defined as the operation of writing a partition depth to a
permanent record.  A permanent record, by definition, is a function
$R : \Nat \to \Real_{\geq 0}$ with $R(t+1) \geq R(t)$ for all $t$ (the
record accumulates, never deletes).  Hence $M(t) = |R^{-1}((0,\infty)) \cap
[0,t]|$ is non-decreasing.  The non-return property follows: no finite
process can reduce $M$, because doing so would require deleting a committed
record entry, which by definition is impossible.  The investigation therefore
has a strict arrow of time given by the accumulation of irreducible residues.
\end{proof}

\begin{remark}
The non-return theorem is the measure-theoretic analogue of the
second law of thermodynamics: the committed information is a non-decreasing
functional.  Each committed individuation increases the count; no process
reverses it.  This gives the contact map investigation a directional
structure independent of any external clock.
\end{remark}

% ==========================================================================
\section{The Minimum Sufficient Contact Map}
\label{sec:mscm}
% ==========================================================================

\subsection{S-Entropy Coordinates}

\begin{definition}[S-Entropy of a Component]
\label{def:s-entropy}
Let $\Part \in \mathfrak{P}(\Omega,d,\mu)$ and $A \in \Part$.  Fix a
$p$-dimensional feature map $\phi : \Omega \to [0,1]^p$.  The
\emph{S-entropy coordinate vector} of $A$ is
\begin{equation}
  \Scoord(A) \;=\; \left( H_\phi^{(1)}(A), \ldots, H_\phi^{(p)}(A) \right)
    \;\in\; [0,1]^p,
\end{equation}
where
\begin{equation}
  H_\phi^{(j)}(A) \;=\; -\int_A \phi_j(x) \log_2 \phi_j(x)\, d\mu(x)
    \;\bigg/\; \mu(A)
\end{equation}
is the normalised Shannon entropy of the $j$-th feature within $A$
\citep{shannon1948mathematical, cover2006elements}.  The \emph{S-entropy
distance} between two components $A_i, A_j \in \Part$ is
\begin{equation}
  \dS(A_i, A_j) \;=\; \left\|\Scoord(A_i) - \Scoord(A_j)\right\|_2.
\end{equation}
\end{definition}

\begin{remark}
The S-entropy coordinates capture the distributional geometry of each
component in feature space.  Components with identical S-entropy coordinates
are, from the perspective of the feature map $\phi$, indistinguishable.
The S-entropy distance is a pseudo-metric on $\Part$; it becomes a metric
when $\phi$ is sufficiently discriminative.
\end{remark}

\subsection{The Contact Graph and Contact Map}

\begin{definition}[Contact Graph]
\label{def:contact-graph}
The \emph{contact graph} of $\Part$ is the undirected graph
$\GC = (\Part, E_C)$ where
\begin{equation}
  E_C = \left\{\{A_i, A_j\} : A_j \in \Sep(A_i, \Part)\right\}.
\end{equation}
Two components are adjacent in $\GC$ if and only if their closures share a
boundary point.
\end{definition}

\begin{definition}[Contact Map]
\label{def:contact-map}
The \emph{contact map} of $(\Part, \phi)$ is the edge-weighting
\begin{equation}
  \CM : E_C \;\longrightarrow\; \Real_{\geq 0}, \qquad
  \CM(A_i, A_j) \;=\; \dS(S(A_i), S(A_j)).
\end{equation}
\end{definition}

\subsection{The Minimum Sufficient Contact Map}

\begin{definition}[Contact Locus]
\label{def:contact-locus}
The \emph{contact locus} of $\Part$ is the sub-collection
\begin{equation}
  \Part_{\bmin} \;=\; \left\{ A \in \Part : \bP(A) = \bmin \right\}
\end{equation}
of components that sit precisely at the partition floor.  The
\emph{minimum sufficient contact map} is the restriction
\begin{equation}
  \CMbmin(A_i, A_j) \;=\; \dS\!\left(\Sbmin(A_i),\, \Sbmin(A_j)\right)
\end{equation}
for $(A_i, A_j) \in E_C$ with $A_i, A_j \in \Part_{\bmin}$, where
$\Sbmin(A) = \Scoord(A)|_{\Part_{\bmin}}$ are the S-entropy coordinates
evaluated at the contact locus.
\end{definition}

\begin{definition}[Partition Depth Inflation]
\label{def:inflation}
A contact map $\CM'$ is a \emph{partition depth inflation} of $\CMbmin$ if
there exists $\delta L \geq 0$ (a non-negative functional) such that
\begin{equation}
  \CM'(A_i, A_j) \;=\; \CMbmin(A_i, A_j) + \delta L(A_i, A_j)
\end{equation}
for all $(A_i, A_j) \in E_C$.
\end{definition}

\subsection{The Minimum Sufficient Lagrangian and Action}

\begin{definition}[Partition Lagrangian]
\label{def:lagrangian}
Given a path $\gamma : [0,T] \to \Omega$ in the bounded resolvable space,
the \emph{partition Lagrangian} associated to a contact map $\CM$ is
\begin{equation}
  L(\gamma, \dot\gamma, t) \;=\; \tfrac{1}{2}\,\mu\!\left(B(\gamma(t), \|\dot\gamma(t)\|)\right)
    \;-\; M_\CM(\gamma(t)),
\end{equation}
where $M_\CM : \Omega \to \Real$ is the \emph{contact potential} induced by
$\CM$:
\begin{equation}
  M_\CM(x) \;=\; \sum_{(A_i,A_j) \in E_C} \CM(A_i,A_j)\,
    \mathbf{1}_{A_i}(x)\,\mathbf{1}_{A_j}(x).
\end{equation}
The \emph{minimum sufficient Lagrangian} is $\Lmin = L|_{\CM = \CMbmin}$.
\end{definition}

\begin{definition}[Minimum Sufficient Action]
\label{def:action}
The \emph{minimum sufficient action} is
\begin{equation}
  \Smin \;=\; \int_0^T \Lmin(\gamma, \dot\gamma, t)\, dt
    \;=\; \int_0^T \left[\tfrac{1}{2}\mu\!\left(B(\gamma,\|\dot\gamma\|)\right)
    - M_{\bmin}(\gamma)\right] dt,
\end{equation}
where $M_{\bmin}$ is the contact potential of $\CMbmin$.  The
\emph{principle of minimum sufficient action} states that the equilibrium
path of a contact map investigation minimises $\Smin$ subject to boundary
conditions $\gamma(0), \gamma(T)$.
\end{definition}

\subsection{The Minimum Sufficiency Theorem}

\begin{theorem}[Minimum Sufficiency]
\label{thm:minimum-sufficiency}
The following four objects are in canonical bijective correspondence:
\begin{enumerate}[label=(\Roman*)]
  \item The partition floor $\bmin$;
  \item The minimum sufficient contact map $\CMbmin$;
  \item The minimum sufficient Lagrangian $\Lmin$;
  \item The minimum sufficient action $\Smin$.
\end{enumerate}
Every other configuration of $\BRS$ is a partition depth inflation of
$\CMbmin$: the general contact map satisfies $\CM = \CMbmin + \delta L$
with $\delta L \geq 0$ pointwise, and the general action satisfies $S =
\Smin + \delta S$ with $\delta S \geq 0$.
\end{theorem}

\begin{proof}
We establish the bijections by constructing explicit maps in each direction
and verifying they are inverses.

\textbf{(I) $\Rightarrow$ (II).}  Given $\bmin$, the contact locus
$\Part_{\bmin}$ is determined by Definition~\ref{def:contact-locus}.  The
S-entropy coordinates $\Sbmin(A)$ for $A \in \Part_{\bmin}$ are determined
by the measure $\mu$ and the feature map $\phi$ (fixed by the space).  Hence
$\CMbmin$ is uniquely determined by $\bmin$.

\textbf{(II) $\Rightarrow$ (I).}  Given $\CMbmin$, the contact locus
$\Part_{\bmin}$ is its domain.  By Definition~\ref{def:partition-depth},
$\bmin = \inf_{A \in \Part_{\bmin}} \bP(A)$.  Since $\CMbmin$ is defined
precisely on the contact locus, $\bmin$ is determined.

\textbf{(II) $\Rightarrow$ (III).}  Given $\CMbmin$, the contact potential
$M_{\bmin}$ is determined by Definition~\ref{def:lagrangian}, and hence
$\Lmin$ is uniquely determined.

\textbf{(III) $\Rightarrow$ (II).}  Given $\Lmin$, the contact potential
$M_{\bmin}$ is the negation of the non-kinetic term of $\Lmin$.  The
contact map values $\CMbmin(A_i,A_j)$ are the coefficients in $M_{\bmin}$,
hence determined.

\textbf{(III) $\Rightarrow$ (IV) and (IV) $\Rightarrow$ (III).}  The action
is a functional of the Lagrangian via time integration; given $\Lmin$, $\Smin$
is determined, and conversely $\Lmin$ is recoverable from $\Smin$ via the
Euler-Lagrange equations \citep{goldstein2002classical, arnold1989mathematical}.

For the inflation claim: let $\CM = \CMbmin + \delta L$ with $\delta L \geq
0$.  Then $M_\CM = M_{\bmin} + M_{\delta L}$ with $M_{\delta L} \geq 0$,
so $L = \Lmin - M_{\delta L} \leq \Lmin$ --- but the action integral is
minimised at $\Smin$, so every inflation increases the action.  No partition
depth inflation can reduce $\CMbmin$, since by definition $\CMbmin$ is already
the minimum.
\end{proof}

\begin{corollary}[Contact Ground State]
\label{cor:ground-state}
The minimum sufficient contact map $\CMbmin$ is the \emph{contact ground
state}: the minimum structure required to establish that every component is
surely not any other component, against the whole system.  Every additional
structural annotation (type declarations, test suites, runtime traces,
lint rules) is a partition depth inflation $\delta L \geq 0$ of this ground
state.  None can reduce the floor to zero; all are upward perturbations.
\end{corollary}

\subsection{The Water and Beads Formalisation}

We formalise the intuition that ``rearranging observable states does not
change the irreducible floor.''

\begin{definition}[Observable States and Irreducible Residue]
\label{def:beads-water}
Let $\Sigma = \{\sigma_1, \ldots, \sigma_N\}$ be a finite set of
\emph{observable states} (test outcomes, type annotations, runtime events).
An \emph{arrangement} is a function $\alpha : \Part \to 2^\Sigma$ assigning
states to components.  The \emph{bead functional} is
\begin{equation}
  \mathcal{B}(\alpha) \;=\; \sum_{A \in \Part} |\alpha(A)|,
\end{equation}
the total count of observable states.  The \emph{water level} is $\bmin$.
\end{definition}

\begin{theorem}[Water Level Invariance]
\label{thm:water-beads}
For any arrangement $\alpha$ and any rearrangement $\alpha' = \pi \circ \alpha$
(where $\pi : 2^\Sigma \to 2^\Sigma$ is any permutation of state assignments),
\begin{equation}
  \bmin(\alpha) \;=\; \bmin(\alpha') \;=\; \bmin,
\end{equation}
i.e., the partition floor is invariant under all rearrangements of observable
states.
\end{theorem}

\begin{proof}
The partition floor $\bmin$ is a property of the measure space $\BRS$ and
the partition $\Part$; it does not depend on the assignment $\alpha$.
Formally: $\bmin = \inf_\Part \inf_{A \in \Part} \inf_{B \in \Sep(A,\Part)}
\mu(B)$.  The measure $\mu$ is a fixed Borel measure on $\Omega$; neither
$\alpha$ nor $\pi$ acts on $\mu$ or on the topological structure of $\Omega$.
Hence $\bmin(\alpha) = \bmin(\alpha') = \bmin$ regardless of the arrangement.
\end{proof}

\begin{remark}
This theorem formalises the water/beads intuition precisely.  The observable
states $\Sigma$ (test results, type errors, coverage reports) are the
``beads'': their arrangement and rearrangement does not change the
``water level'' $\bmin$.  A system with $N$ passing tests and a system with
$N$ failing tests have the same partition floor if they have the same
underlying bounded resolvable space.  The contact map generator's task is
to measure $\bmin$ accurately --- not to remove it.
\end{remark}

% ==========================================================================
\section{Software Systems as Bounded Resolvable Spaces}
\label{sec:software}
% ==========================================================================

\subsection{The System Graph}

\begin{definition}[Software System]
\label{def:software-system}
A \emph{software system} is a quadruple $\mathfrak{S} = (U, E, A, \Pi)$
where:
\begin{enumerate}[label=(\roman*)]
  \item $U = \{u_1, \ldots, u_N\}$ is a finite set of \emph{units} (functions,
    classes, modules, services, or any other named computation boundary);
  \item $E \subseteq U \times U$ is a set of directed \emph{edges}
    representing invocation, data-flow, or dependency relations;
  \item $A : U \to 2^\mathcal{T}$ assigns to each unit an \emph{aperture set}
    --- the collection of types it consumes;
  \item $\Pi : U \to 2^\mathcal{T}$ assigns to each unit a
    \emph{production set} --- the collection of types it produces;
\end{enumerate}
where $\mathcal{T}$ is a fixed type universe.  The \emph{system graph}
$\GS = (U, E)$ is the underlying directed graph.
\end{definition}

\begin{definition}[Software Metric Space]
\label{def:software-metric}
Given $\mathfrak{S}$, define a metric on $U$ by
\begin{equation}
  d_\mathfrak{S}(u_i, u_j) \;=\;
    1 - \frac{|A(u_i) \cap A(u_j)| + |\Pi(u_i) \cap \Pi(u_j)|}
             {|A(u_i) \cup A(u_j)| + |\Pi(u_i) \cup \Pi(u_j)|},
\end{equation}
the complement of the Jaccard similarity of the combined aperture/production
type sets.  Define a measure on $U$ by $\mu(u_i) = w_i$ where $w_i$ is the
normalised coupling weight (e.g., proportional to in-degree in $\GS$).
\end{definition}

\begin{proposition}
$(U, d_\mathfrak{S}, \mu)$ is a bounded resolvable space (in the finite,
discrete sense) with resolution data $(1, w_{\min}, 1)$, where $w_{\min} =
\min_i w_i > 0$.
\end{proposition}

\begin{proof}
$U$ is finite, hence compact.  $\mu$ is a finite measure with $\mu(u_i) =
w_i > 0$ for all $i$ (by assumption of positive coupling weight).  The
resolution floor holds trivially: $\mu(B(u_i, r)) \geq w_i \geq w_{\min}$
for $r \geq 0$.
\end{proof}

\subsection{Holonomy of System Cycles}

\begin{definition}[Cycle Holonomy]
\label{def:holonomy}
Let $c = (u_{i_1}, u_{i_2}, \ldots, u_{i_k}, u_{i_1})$ be a directed cycle
in $\GS$.  The \emph{specification type flow} along $c$ is the sequence of
type transitions
\begin{equation}
  \tau(c) \;=\; \Pi(u_{i_1}) \to A(u_{i_2}) \to \Pi(u_{i_2}) \to \cdots
    \to A(u_{i_1}).
\end{equation}
The \emph{holonomy} of $c$ is
\begin{equation}
  \hol(c) \;=\; \left\| \tau(c) - \tau^{\mathrm{spec}}(c) \right\|,
\end{equation}
where $\tau^{\mathrm{spec}}(c)$ is the specified type flow and $\|\cdot\|$
is a norm on the space of type-transition sequences.
\end{definition}

\begin{remark}
Holonomy measures how much the actual type flow around a cycle deviates from
its specification.  A system with $\hol(c) = 0$ for all cycles is
\emph{specificationally consistent} around all loops.  Non-zero holonomy
indicates a coordination failure: the types produced by traversing the cycle
do not match the types expected.
\end{remark}

\subsection{S-Entropy of Software Units}

\begin{definition}[S-Entropy of a Software Unit]
\label{def:s-entropy-software}
For a software unit $u \in U$, fix a $p$-dimensional feature map
$\phi_\mathfrak{S} : U \to [0,1]^p$ encoding structural properties
(e.g.\ cyclomatic complexity, coupling degree, type-set cardinality,
dependency depth).  The \emph{S-entropy coordinate vector} of $u$ is
\begin{equation}
  \Scoord(u) \;=\; \left(-\phi_j(u)\log_2\phi_j(u)\right)_{j=1}^p
    \;\in\; [0,1]^p.
\end{equation}
The \emph{S-entropy} of $u$ is $S(u) = \|\Scoord(u)\|_1 \in [0, p\log_2 e]$.
The floor $S_{\flat} = \min_{u \in U} S(u)$ is strictly positive for any
non-trivially structured system.
\end{definition}

\subsection{Five Coordination Regimes}

\begin{definition}[Static Order Parameter]
\label{def:order-parameter}
The \emph{static order parameter} of $\mathfrak{S}$ is
\begin{equation}
  \Rest \;=\; \exp\!\left(-\,\overline{\tau}\right),
\end{equation}
where $\overline{\tau}$ is the mean tension:
\begin{equation}
  \overline{\tau} \;=\; \frac{1}{|E_C|} \sum_{(u_i, u_j) \in E_C}
    \dS(S(u_i), S(u_j)) \cdot \frac{\hol(c_{ij})}{\hol_{\max} + \varepsilon},
\end{equation}
with $c_{ij}$ the shortest cycle containing the edge $(u_i,u_j)$,
$\hol_{\max} = \max_c \hol(c)$, and $\varepsilon > 0$ a regulariser.
\end{definition}

\begin{definition}[Coordination Regimes]
\label{def:coordination-regimes}
The coordination regimes of a software system are determined by $\Rest$:

\begin{center}
\begin{tabular}{lll}
  \toprule
  Regime & $\Rest$ range & Characterisation \\
  \midrule
  Turbulent & $[0, 0.3)$
    & Maximum partition inflation; near-zero contact coherence \\
  Aperture-dominated & $[0.3, 0.5)$
    & Type boundaries dominate; structural coupling weak \\
  Hierarchical cascade & $[0.5, 0.8)$
    & Layered dependency structure; partial contact coherence \\
  Coherent & $[0.8, 0.95)$
    & High contact coherence; near minimum sufficient structure \\
  Phase-locked & $[0.95, 1]$
    & Near contact ground state; $\CM \approx \CMbmin$ \\
  \bottomrule
\end{tabular}
\end{center}
\end{definition}

\begin{remark}
$\Rest$ is an order parameter in the statistical mechanics sense: it measures
proximity of the system's actual contact map to the minimum sufficient contact
map.  $\Rest \to 1$ corresponds to $\CM \to \CMbmin$; $\Rest \to 0$
corresponds to maximally inflated partition depth ($\CM \gg \CMbmin$).
\end{remark}

\subsection{Local Invisibility Theorem}

\begin{theorem}[Local Invisibility]
\label{thm:local-invisibility}
Let $\mathfrak{S} = (U, E, A, \Pi)$ and let $\val_u : \mathcal{P}_u \to
\{0,1\}$ be a local verification predicate for unit $u$, where $\mathcal{P}_u$
is the set of properties of $u$ in isolation.  Then:
\begin{equation}
  \forall u \in U : \val_u(s) = 1 \;\not\Rightarrow\;
    \CM_\mathfrak{S} = \CMbmin^\mathfrak{S}.
\end{equation}
That is, universal local validity does not imply that the system's contact
map is at the minimum sufficient ground state.
\end{theorem}

\begin{proof}
Construct a counterexample.  Let $U = \{u_1, u_2, u_3\}$ with edges
$E = \{(u_1,u_2),(u_2,u_3),(u_3,u_1)\}$ (a 3-cycle).  Define aperture and
production sets so that each $u_i$ satisfies its local type contract in
isolation: $\Pi(u_i) \subseteq A(u_{i+1 \mod 3})$.  Now perturb the types
along the cycle by inserting a coercion $\kappa : T_1 \to T_1'$ that is
invisible to each local verifier (e.g., a silent widening conversion).  Each
$\val_{u_i}(s) = 1$ because the local type contract is satisfied.  However,
the holonomy of the cycle is $\hol(c) = \|\kappa\| > 0$, so the contact map
$\CM_\mathfrak{S}(u_i, u_j) > \CMbmin^\mathfrak{S}(u_i, u_j)$ for some
$(u_i,u_j)$.  Hence global contact map coherence fails despite universal
local validity.

The general argument follows from Rice's theorem \citep{rice1953classes}:
the global semantic property $\CM_\mathfrak{S} = \CMbmin$ is non-trivial and
hence undecidable from local information alone; local predicates $\val_u$
capture only decidable local properties and are therefore insufficient.
\end{proof}

% ==========================================================================
\section{The Contact Map Generator}
\label{sec:generator}
% ==========================================================================

\subsection{What the Contact Map Generator Is}

The \emph{Contact Map Generator} (CMG) is a system that, given a software
system $\mathfrak{S}$, produces a sequence of progressively refined estimates
\begin{equation}
  \hat{\CM}^{(0)} \;\leq\; \hat{\CM}^{(1)} \;\leq\; \cdots \;\leq\; \hat{\CM}^{(T)}
    \;\xrightarrow{T\to\infty}\; \CMbmin^\mathfrak{S}
\end{equation}
of the minimum sufficient contact map, where $\leq$ denotes refinement in the
sense of decreasing partition depth inflation ($\delta L$ monotone
non-increasing).

The CMG is not a test runner.  It is not a linter.  It is not a type checker.
It is a \emph{partition machine}: it accepts any process --- static analysis,
dynamic tracing, LLM-based semantic inference, user-provided structural
intuition --- as a \emph{partition operation} that deposits information about
$\CMbmin^\mathfrak{S}$.  Each operation contributes evidence; the CMG
synthesises all evidence into a running estimate.

\subsection{Partition Operations}

\begin{definition}[Partition Operation]
\label{def:partition-operation}
A \emph{partition operation} on $\mathfrak{S}$ is a computable function
\begin{equation}
  \mathcal{O} : \mathfrak{S} \;\longrightarrow\; \hat{\CM}
\end{equation}
producing a partial contact map estimate $\hat{\CM}$ (a partial function on
$E_C$).  Partition operations include:
\begin{enumerate}[label=(\alph*)]
  \item \textbf{Static}: call graph extraction, type aperture analysis,
    dependency edge analysis \citep{nielson1999principles, aho2006compilers};
  \item \textbf{Dynamic}: runtime trace analysis, coverage profiling,
    execution-path sampling;
  \item \textbf{Semantic}: LLM-based inference of coordination boundaries
    from code text;
  \item \textbf{Structural}: user-provided assertions about architectural
    boundaries.
\end{enumerate}
\end{definition}

\begin{definition}[Composition of Partition Operations]
\label{def:composition}
Given two partition operations $\mathcal{O}_1, \mathcal{O}_2$ producing
estimates $\hat{\CM}_1, \hat{\CM}_2$, their \emph{composition} is
\begin{equation}
  (\mathcal{O}_1 \oplus \mathcal{O}_2)(\mathfrak{S}) \;=\;
    \hat{\CM}_1 \wedge \hat{\CM}_2,
\end{equation}
where $\wedge$ takes the pointwise minimum (the more conservative estimate
on each edge), respecting the monotone non-inflation property.
\end{definition}

\subsection{The Cessation Theorem}

\begin{theorem}[Cessation]
\label{thm:cessation}
Let $(\mathcal{O}^{(t)})_{t \geq 0}$ be a sequence of partition operations
applied by the CMG, with marginal returns $\Delta_t = \|\hat{\CM}^{(t)} -
\hat{\CM}^{(t-1)}\|$.  Then $\Delta_t \to \bmin > 0$ as $t \to \infty$;
the CMG halts (ceases further operations) when $\Delta_t < \bmin$.
\end{theorem}

\begin{proof}
By Theorem~\ref{thm:floor-positive}, $\bmin > 0$.  Each partition operation
$\mathcal{O}^{(t)}$ deposits at least $\bmin$ information (by
Definition~\ref{def:partition-depth} and the floor positivity), so $\Delta_t
\geq \bmin$ for all $t$ until the estimate converges to $\CMbmin$.  Once
$\hat{\CM}^{(t)} = \CMbmin$ (up to the resolution of the feature map $\phi$),
further operations produce $\Delta_t = 0 < \bmin$ (no new information above
the floor), and the cessation condition is met.  Hence the CMG halts in finite
time --- at or near the contact ground state.
\end{proof}

% ==========================================================================
\section{The Wind Tunnel DSL}
\label{sec:dsl}
% ==========================================================================

\subsection{Overview and Design Principles}

The \emph{Wind Tunnel DSL} is the language in which contact map
investigations are specified and recorded.  It is not a test language and not
a configuration language.  Its design principles are:

\begin{enumerate}[label=(\roman*)]
  \item \textbf{Partition-first}: every construct corresponds to a partition
    operation or a claim about partition depths.
  \item \textbf{Co-authored}: both users and the CMG system write valid DSL.
    The same syntax is used by both.
  \item \textbf{Self-describing}: a completed DSL execution file is a
    complete, reproducible record of a contact map investigation.
  \item \textbf{Monotone}: committed partition claims cannot be undone within
    a script (Non-Return Theorem, Theorem~\ref{thm:non-return}).
  \item \textbf{Floor-aware}: every \texttt{assert} is a minimum individuation
    claim (a lower bound on $\CMbmin$), never a truth-at-a-point claim.
\end{enumerate}

\subsection{Formal Grammar}

The Wind Tunnel DSL is defined by the following context-free grammar in BNF
form.

\begin{definition}[Wind Tunnel Grammar]
\label{def:grammar}
\begin{align*}
\langle\textit{program}\rangle &\;::=\; \langle\textit{block}\rangle^+ \\[4pt]
\langle\textit{block}\rangle &\;::=\;
  \langle\textit{scope-block}\rangle \;\mid\; \langle\textit{analyse-block}\rangle \\
  &\quad\;\mid\; \langle\textit{compose-block}\rangle \;\mid\;
    \langle\textit{assert-block}\rangle \\
  &\quad\;\mid\; \langle\textit{recurse-block}\rangle \;\mid\;
    \langle\textit{report-block}\rangle \\[4pt]
\langle\textit{scope-block}\rangle &\;::=\;
  \texttt{scope}\;\langle\textit{id}\rangle\;\texttt{\{}\;
  \langle\textit{scope-body}\rangle\;\texttt{\}} \\
\langle\textit{scope-body}\rangle &\;::=\;
  \langle\textit{target-decl}\rangle^+\;\langle\textit{filter-decl}\rangle^* \\
\langle\textit{target-decl}\rangle &\;::=\;
  \texttt{target:}\;\langle\textit{glob-pattern}\rangle \\
\langle\textit{filter-decl}\rangle &\;::=\;
  \texttt{exclude:}\;\langle\textit{glob-pattern}\rangle
  \;\mid\; \texttt{depth:}\;\langle\textit{nat}\rangle \\[4pt]
\langle\textit{analyse-block}\rangle &\;::=\;
  \texttt{analyse}\;\langle\textit{method}\rangle\;\texttt{\{}\;
  \langle\textit{analyse-body}\rangle\;\texttt{\}} \\
\langle\textit{method}\rangle &\;::=\;
  \texttt{static} \;\mid\; \texttt{dynamic} \;\mid\;
  \texttt{semantic} \;\mid\; \texttt{hf} \;\mid\; \langle\textit{id}\rangle \\
\langle\textit{analyse-body}\rangle &\;::=\;
  \langle\textit{step}\rangle^+ \\
\langle\textit{step}\rangle &\;::=\;
  \texttt{extract:}\;\langle\textit{id}\rangle
  \;\mid\; \texttt{model:}\;\langle\textit{string}\rangle
  \;\mid\; \texttt{query:}\;\langle\textit{string}\rangle
  \;\mid\; \texttt{method:}\;\langle\textit{id}\rangle \\[4pt]
\langle\textit{compose-block}\rangle &\;::=\;
  \texttt{compose}\;\texttt{\{}\;
  \langle\textit{compose-body}\rangle\;\texttt{\}} \\
\langle\textit{compose-body}\rangle &\;::=\;
  \langle\textit{merge-decl}\rangle^+ \\
\langle\textit{merge-decl}\rangle &\;::=\;
  \texttt{merge:}\;\texttt{[}\;\langle\textit{id-list}\rangle\;\texttt{]}
  \;\mid\; \texttt{resolve:}\;\langle\textit{resolve-strategy}\rangle \\
\langle\textit{resolve-strategy}\rangle &\;::=\;
  \texttt{overlaps by minimum\_depth}
  \;\mid\; \texttt{overlaps by maximum\_depth}
  \;\mid\; \langle\textit{id}\rangle \\[4pt]
\langle\textit{assert-block}\rangle &\;::=\;
  \texttt{assert}\;\langle\textit{assertion-expr}\rangle\;\texttt{\{}\;
  \langle\textit{assert-body}\rangle\;\texttt{\}} \\
\langle\textit{assertion-expr}\rangle &\;::=\;
  \langle\textit{individuation-claim}\rangle
  \;\mid\; \langle\textit{distance-claim}\rangle
  \;\mid\; \langle\textit{regime-claim}\rangle \\
\langle\textit{individuation-claim}\rangle &\;::=\;
  \texttt{individuation} \\
\langle\textit{distance-claim}\rangle &\;::=\;
  \texttt{distance} \\
\langle\textit{regime-claim}\rangle &\;::=\;
  \texttt{regime} \\
\langle\textit{assert-body}\rangle &\;::=\;
  \texttt{claim:}\;\langle\textit{claim-expr}\rangle
  \;\langle\textit{on-fail-decl}\rangle^? \\
\langle\textit{claim-expr}\rangle &\;::=\;
  \texttt{distance(}\langle\textit{id}\rangle\texttt{,}\langle\textit{id}\rangle\texttt{)}\;
  \texttt{>=}\;\langle\textit{float}\rangle \\
  &\quad\;\mid\; \texttt{R\_est}\;\texttt{>=}\;\langle\textit{float}\rangle
  \;\mid\; \langle\textit{id}\rangle\;\langle\textit{op}\rangle\;\langle\textit{float}\rangle \\
\langle\textit{on-fail-decl}\rangle &\;::=\;
  \texttt{on\_fail:}\;\langle\textit{fail-action}\rangle \\
\langle\textit{fail-action}\rangle &\;::=\;
  \texttt{narrow\_scope(}\langle\textit{id}\rangle\texttt{,}\langle\textit{id}\rangle\texttt{)}
  \;\mid\; \langle\textit{recurse-block}\rangle
  \;\mid\; \langle\textit{id}\rangle \\[4pt]
\langle\textit{recurse-block}\rangle &\;::=\;
  \texttt{recurse}\;\langle\textit{condition}\rangle\;\texttt{\{}\;
  \langle\textit{recurse-body}\rangle\;\texttt{\}} \\
\langle\textit{condition}\rangle &\;::=\;
  \texttt{until}\;\langle\textit{threshold-expr}\rangle
  \;\mid\; \texttt{while}\;\langle\textit{cond-expr}\rangle \\
\langle\textit{threshold-expr}\rangle &\;::=\;
  \texttt{marginal\_return <}\;\langle\textit{id}\rangle
  \;\mid\; \langle\textit{id}\rangle\;\langle\textit{op}\rangle\;\langle\textit{float}\rangle \\
\langle\textit{recurse-body}\rangle &\;::=\;
  \langle\textit{block}\rangle^+ \\[4pt]
\langle\textit{report-block}\rangle &\;::=\;
  \texttt{report}\;\texttt{\{}\;\langle\textit{report-body}\rangle\;\texttt{\}} \\
\langle\textit{report-body}\rangle &\;::=\;
  \langle\textit{output-decl}\rangle^+ \\
\langle\textit{output-decl}\rangle &\;::=\;
  \texttt{output:}\;\langle\textit{id}\rangle
  \;\mid\; \texttt{format:}\;\langle\textit{id}\rangle
\end{align*}
\end{definition}

\subsection{Operational Semantics}

We define a small-step operational semantics for the DSL.

\begin{definition}[DSL State]
\label{def:dsl-state}
A \emph{DSL state} is a triple $\Gamma = (\hat{\CM}, \mathcal{C}, \mathcal{R})$
where:
\begin{enumerate}[label=(\roman*)]
  \item $\hat{\CM}$ is the current contact map estimate;
  \item $\mathcal{C} \subseteq E_C$ is the set of committed edges (pairs
    whose partition depth has been certified);
  \item $\mathcal{R} \in \{\textit{pass}, \textit{fail}, \textit{narrow},
    \textit{recurse}\}$ is the current resolution flag.
\end{enumerate}
The initial state is $\Gamma_0 = (\hat{\CM}^{(0)}, \emptyset, \textit{pass})$
where $\hat{\CM}^{(0)}$ is the trivial (all-unknown) estimate.
\end{definition}

\begin{definition}[Block Reduction Rules]
\label{def:reduction-rules}
The reduction relation $\Gamma \xrightarrow{b} \Gamma'$ for each block
type $b$ is defined as follows.

\paragraph{Scope.}
$(\hat{\CM}, \mathcal{C}, \mathcal{R})
  \xrightarrow{\texttt{scope}(id, \textit{targets})}
  (\hat{\CM}, \mathcal{C}, \mathcal{R})[\textit{active-scope} := \textit{targets}]$.
The scope block declares the active partition region.  It does not modify the
contact map or commit any partition.

\paragraph{Analyse.}
$(\hat{\CM}, \mathcal{C}, \mathcal{R})
  \xrightarrow{\texttt{analyse}(\mathit{method}, \textit{steps})}
  (\hat{\CM} \wedge \hat{\CM}_{\mathit{method}}, \mathcal{C}, \mathcal{R})$,
where $\hat{\CM}_{\mathit{method}}$ is the partial contact map produced by
partition operation $\mathcal{O}_{\mathit{method}}$ applied to the active
scope.  The composition $\wedge$ is pointwise minimum.

\paragraph{Compose.}
$(\hat{\CM}, \mathcal{C}, \mathcal{R})
  \xrightarrow{\texttt{compose}(\textit{merge-list})}
  (\bigwedge_j \hat{\CM}_j, \mathcal{C}, \mathcal{R})$,
where $\hat{\CM}_j$ are the estimates named in the merge list and $\bigwedge$
is pointwise minimum over all named estimates.

\paragraph{Assert (success).}
If $\textit{claim}$ evaluates to $\top$ in state $\Gamma$:
$(\hat{\CM}, \mathcal{C}, \mathcal{R})
  \xrightarrow{\texttt{assert}(\textit{claim})}
  (\hat{\CM}, \mathcal{C} \cup \{\textit{certified-edges}\}, \textit{pass})$.
The certified edges (those satisfying the claim) are committed.

\paragraph{Assert (failure).}
If $\textit{claim}$ evaluates to $\bot$:
$(\hat{\CM}, \mathcal{C}, \mathcal{R})
  \xrightarrow{\texttt{assert}(\textit{claim},\textit{on-fail}=f)}
  (\hat{\CM}, \mathcal{C}, \textit{narrow})[f]$,
where $[f]$ indicates that the fail-action $f$ is queued for execution.
The block does not halt; it narrows the scope and continues.

\paragraph{Recurse.}
$(\hat{\CM}, \mathcal{C}, \mathcal{R})
  \xrightarrow{\texttt{recurse}(\textit{cond}, \textit{body})}
  (\hat{\CM}^*, \mathcal{C}^*, \mathcal{R}^*)$,
where $(\hat{\CM}^*, \mathcal{C}^*, \mathcal{R}^*)$ is the fixed point of
repeatedly applying $\textit{body}$ until $\textit{cond}$ is satisfied.  The
termination condition is guaranteed by the Cessation Theorem
(Theorem~\ref{thm:cessation}).

\paragraph{Report.}
$(\hat{\CM}, \mathcal{C}, \mathcal{R})
  \xrightarrow{\texttt{report}(\textit{outputs})}
  (\hat{\CM}, \mathcal{C}, \mathcal{R})$
with side effect: emit the named outputs (contact map, regime classification,
inflation above floor) to the specified format.  The state is unchanged.
\end{definition}

\subsection{Semantic Properties}

\begin{theorem}[Compositional Soundness]
\label{thm:compositional}
Let $b_1, b_2$ be successive blocks in a DSL program with reduction
$\Gamma \xrightarrow{b_1} \Gamma' \xrightarrow{b_2} \Gamma''$.  Then:
\begin{enumerate}[label=(\roman*)]
  \item \textbf{Monotone commitment}: $\mathcal{C} \subseteq \mathcal{C}'
    \subseteq \mathcal{C}''$ (committed edges are never revoked).
  \item \textbf{Non-inflation}: $\hat{\CM}'' \leq \hat{\CM}' \leq \hat{\CM}$
    pointwise (estimates are monotone non-increasing in partition depth).
  \item \textbf{Floor preservation}: $\hat{\CM}'' \geq \CMbmin$ pointwise
    (no estimate can go below the floor).
\end{enumerate}
\end{theorem}

\begin{proof}
(i) Follows from the Assert rule: commitment extends $\mathcal{C}$ but never
removes from it.  The Non-Return Theorem (Theorem~\ref{thm:non-return})
applies.

(ii) Each Analyse and Compose step applies $\wedge$ (pointwise minimum),
so the estimate can only decrease or remain constant.

(iii) Every partition operation $\mathcal{O}$ produces estimates bounded
below by $\bmin$ (Theorem~\ref{thm:floor-positive}), so $\hat{\CM}_\mathcal{O}
\geq \CMbmin$.  The $\wedge$ of estimates bounded below by $\CMbmin$ is itself
bounded below by $\CMbmin$.
\end{proof}

\subsection{Execution Modes}

\begin{enumerate}[label=\textbf{Mode \arabic*.}, leftmargin=*]
\item \textbf{User-only.}  The user writes a complete DSL program including
  all \texttt{scope}, \texttt{analyse}, \texttt{compose}, \texttt{assert},
  and \texttt{report} blocks.  The CMG executes the program faithfully,
  performing the specified partition operations and returning the contact map
  as specified.  The system does not add any blocks.

\item \textbf{System-only.}  The user provides only a \texttt{scope} block
  (or a target path).  The CMG writes and executes the remainder of the
  program --- selecting methods, writing \texttt{analyse} and \texttt{compose}
  blocks, generating \texttt{assert} claims based on its current best estimate
  of $\CMbmin^\mathfrak{S}$, and recursing until the Cessation Theorem applies.

\item \textbf{Collaborative.}  The user writes \texttt{scope} and
  \texttt{assert} blocks (encoding architectural intuition and minimum
  individuation claims).  The CMG fills in the \texttt{analyse}, \texttt{compose},
  and \texttt{recurse} blocks.  Discrepancies between the user's assert claims
  and the CMG's computed contact map are surfaced as new \texttt{report} blocks
  indicating where the user's structural model diverges from the measured
  contact map.
\end{enumerate}

\subsection{Worked Examples}

\subsubsection{Example 1: Authentication Module Analysis}

The following script investigates the contact structure of an authentication
module.  It is a collaborative script: the user wrote \texttt{scope} and
\texttt{assert}; the CMG wrote \texttt{analyse}, \texttt{compose}, and the
\texttt{recurse} body.

\begin{lstlisting}[caption={Authentication module contact map investigation},
                   label={lst:auth}]
# Contact map investigation: auth module
# User-authored: scope, assert blocks
# System-authored: analyse, compose, recurse

scope authentication_module {
  target: "src/auth/**/*.ts"
  exclude: "**/*.test.ts"
  exclude: "**/*.d.ts"
  depth: 3
}

# System: extract structural partition operations
analyse static {
  extract: call_graph
  extract: type_apertures
  extract: dependency_edges
  extract: export_surface
}

# System: semantic inference of coordination boundaries
analyse semantic {
  method: hf
  model: "bigcode/starcoder2-7b"
  query: "identify coordination boundaries and
          interface contracts between components"
}

# System: compose static and semantic estimates
compose {
  merge: [static.call_graph,
          static.type_apertures,
          semantic.boundaries]
  resolve: overlaps by minimum_depth
}

# User: architectural individuation claim
# (TokenValidator and UserStore should be
#  at least 0.6 apart in S-entropy space)
assert individuation {
  claim: distance(auth.TokenValidator,
                  auth.UserStore) >= 0.6
  on_fail: narrow_scope(auth.TokenValidator,
                        auth.UserStore)
}

# User: regime claim (coherent or better)
assert regime {
  claim: R_est >= 0.8
  on_fail: recurse until marginal_return < floor {
    analyse dynamic {
      extract: runtime_call_trace
      extract: type_coercions
    }
    compose {
      merge: [current, dynamic.call_trace]
      resolve: overlaps by minimum_depth
    }
  }
}

# System: assert holonomy of all 3-cycles
assert individuation {
  claim: distance(auth.SessionManager,
                  auth.TokenValidator) >= 0.45
  on_fail: narrow_scope(auth.SessionManager,
                        auth.TokenValidator)
}

assert individuation {
  claim: distance(auth.UserStore,
                  auth.PasswordHasher) >= 0.55
  on_fail: narrow_scope(auth.UserStore,
                        auth.PasswordHasher)
}

report {
  output: contact_map
  output: regime_classification
  output: inflation_above_floor
  output: holonomy_per_cycle
  format: json
}
\end{lstlisting}

\subsubsection{Example 2: System-Only Investigation}

Here the user provides only a scope; the CMG generates the complete
investigation script.

\begin{lstlisting}[caption={System-generated contact map investigation},
                   label={lst:system-only}]
# Contact map investigation: payment service
# User-authored: scope only
# System-authored: all remaining blocks

scope payment_service {
  target: "services/payment/**"
  exclude: "**/vendor/**"
}

# System: layered partition analysis
analyse static {
  extract: call_graph
  extract: type_apertures
  extract: dependency_edges
  extract: api_surface
  extract: shared_state
}

analyse semantic {
  method: hf
  model: "bigcode/starcoder2-7b"
  query: "identify domain boundaries,
          transaction boundaries,
          and external service interfaces"
}

analyse static {
  method: holonomy
  extract: cycle_detection
  extract: holonomy_per_cycle
}

compose {
  merge: [static.call_graph,
          static.type_apertures,
          static.shared_state,
          semantic.domain_boundaries,
          static.holonomy_per_cycle]
  resolve: overlaps by minimum_depth
}

# System: floor-aware regime check
assert regime {
  claim: R_est >= 0.5
  on_fail: recurse until marginal_return < floor {
    analyse semantic {
      method: hf
      model: "bigcode/starcoder2-7b"
      query: "identify type coercions
              and implicit contracts"
    }
    compose {
      merge: [current,
              semantic.type_coercions]
      resolve: overlaps by minimum_depth
    }
  }
}

# System: assert minimum individuation for
# highest-tension edge pairs
assert individuation {
  claim: distance(payment.Processor,
                  payment.Ledger) >= floor
  on_fail: narrow_scope(payment.Processor,
                        payment.Ledger)
}

assert individuation {
  claim: distance(payment.Gateway,
                  payment.Validator) >= floor
  on_fail: narrow_scope(payment.Gateway,
                        payment.Validator)
}

report {
  output: contact_map
  output: regime_classification
  output: inflation_above_floor
  output: tension_per_edge
  output: holonomy_summary
  format: json
}
\end{lstlisting}

\subsubsection{Example 3: Discrepancy Surfacing}

This example illustrates a collaborative investigation where the user's
structural intuition (assert) is inconsistent with the measured contact map.
The CMG surfaces the discrepancy as a new report block.

\begin{lstlisting}[caption={Discrepancy between user model and measured contact map},
                   label={lst:discrepancy}]
# User believes the API layer and domain layer
# are well-separated (distance >= 0.8).
# CMG finds actual distance is 0.41.

scope api_domain_boundary {
  target: "src/**"
  exclude: "src/infrastructure/**"
}

analyse static {
  extract: call_graph
  extract: type_apertures
}

analyse semantic {
  method: hf
  model: "bigcode/starcoder2-7b"
  query: "characterise coupling between
          API handlers and domain objects"
}

compose {
  merge: [static.call_graph,
          semantic.coupling]
  resolve: overlaps by minimum_depth
}

# User assertion: architectural intent
assert individuation {
  claim: distance(api.handlers,
                  domain.aggregates) >= 0.8
  on_fail: narrow_scope(api.handlers,
                        domain.aggregates)
}

# CMG report: assertion failed.
# Measured distance: 0.41
# Regime: aperture-dominated (R_est = 0.38)
# The api.handlers layer directly references
# 14 domain.aggregates types, producing
# high S-entropy overlap.
# Inflation above floor: 1.93x

report {
  output: contact_map
  output: regime_classification
  output: inflation_above_floor
  output: discrepancy_report
  format: json
}
\end{lstlisting}

% ==========================================================================
\section{Discussion}
\label{sec:discussion}
% ==========================================================================

\subsection{Relation to Classical Testing}

Classical testing \citep{clarke1999model, baier2008principles} evaluates
predicates at points in a system's execution space.  Model checking
exhaustively verifies all paths through a finite state space; property-based
testing samples from a distribution of inputs.  Both paradigms are
\emph{point-evaluation} paradigms: the unit of information is a Boolean value
at a specific execution point.

The contact map framework is not a replacement for these paradigms within
their own domain of applicability.  It is a different level of analysis.  A
test failure is \emph{evidence} that some pair of components is more similar
(in S-entropy space) than the minimum sufficient contact map would predict ---
that is, a test failure is a signal about the contact map, not a refutation of
a system-level property.

More precisely: let $\phi(x, f(x)) = 0$ be a test failure at input $x$.  This
event carries information $-\log_2 p_{\text{fail}}$ bits about the execution,
where $p_{\text{fail}}$ is the prior probability of failure.  In the contact
map framework, this information is a partition operation of type ``dynamic''
that updates the estimate $\hat{\CM}$.  The test failure does not directly
certify a partition depth; it contributes evidence that the contact distance
$\CM(u_i, u_j)$ between the failing unit $u_i$ and the unit responsible for
the violated property $u_j$ is less than previously estimated.

\subsection{The First Cannot Answer the Second}

The central methodological claim of this paper is categorical:

\begin{proposition}
No finite collection of point-evaluation predicates can determine $\CMbmin$ for
a non-trivial bounded resolvable space.
\end{proposition}

\begin{proof}
$\CMbmin$ is defined over the contact locus $\Part_{\bmin}$, which is an
uncountably-indexed family of partition configurations (parameterised by all
valid partitions $\Part \in \mathfrak{P}$).  A finite collection of
point-evaluation predicates can certify at most countably many bits of
information \citep{cover2006elements}, which is insufficient to determine an
uncountable family.  Moreover, by Rice's theorem \citep{rice1953classes}, the
semantic property defining $\CMbmin$ is non-trivial and undecidable; hence no
computable point-evaluation oracle can determine it.
\end{proof}

\subsection{The Second Subsumes the First}

\begin{proposition}
Every classical test assertion is a partition operation in the contact map
framework (specifically, a dynamic-type analyse block with boolean output).
\end{proposition}

\begin{proof}
A test assertion $\phi(x_i, f(x_i)) \in \{0,1\}$ can be encoded as a
partition operation $\mathcal{O}_{\phi,x_i} : \mathfrak{S} \to \hat{\CM}$
that updates the contact map estimate based on the pass/fail outcome.  A
passing test increases confidence that $\CM(u_i, u_j) \geq \theta$ for some
threshold $\theta$; a failing test provides evidence that $\CM(u_i, u_j) <
\theta$.  Hence every classical test is a special case of the partition
operations composable within the DSL's \texttt{analyse dynamic} block.
\end{proof}

\subsection{Falsifiability and Scope}

The framework is falsifiable in the following sense.

\begin{proposition}
The claim ``$\CMbmin^\mathfrak{S}(\cdot, \cdot) \geq \theta$ for all adjacent
pairs'' is falsified by exhibiting a pair $(u_i, u_j) \in E_C$ with
$\dS(S(u_i), S(u_j)) < \theta$, which is detectable by any partition
operation that discriminates $u_i$ from $u_j$.
\end{proposition}

The framework's scope is limited to systems that can be modelled as bounded
resolvable spaces --- i.e., finite software systems with well-defined unit
boundaries and measurable coupling weights.  It does not apply to systems with
zero-measure components (vacuous units) or systems with no measurable
separation (fully monolithic systems with no partition structure).

\subsection{The DSL as Audit Trail}

A completed DSL execution file serves as a complete audit trail for a contact
map investigation.  By the Non-Return Theorem (Theorem~\ref{thm:non-return})
and Compositional Soundness (Theorem~\ref{thm:compositional}), the sequence of
committed partitions in $\mathcal{C}$ is a monotone non-decreasing record of
the investigation's findings.  This has practical implications:

\begin{enumerate}[label=(\roman*)]
  \item \textbf{Reproducibility}: the script can be re-executed to reproduce
    the contact map estimate (modulo non-determinism in LLM inference);
  \item \textbf{Auditability}: the committed partitions form a verifiable
    record of what was individuated and when;
  \item \textbf{Extensibility}: new \texttt{analyse} blocks can be appended
    to refine the estimate without invalidating prior committed partitions.
\end{enumerate}

% ==========================================================================
\section{Conclusion}
\label{sec:conclusion}
% ==========================================================================

We have introduced a measure-theoretic framework for software system
characterisation grounded in the concept of the \emph{minimum sufficient
contact map}.  The framework rests on three main results.

First, the \emph{Partition Floor Theorem} (Theorem~\ref{thm:floor-positive}),
proved by three independent methods (geometric, representational, economic),
establishes that no investigation strategy can drive the separator between
adjacent components to zero measure.  The irreducible floor $\bmin > 0$ is
a property of the space, not of the investigation method.

Second, the \emph{Minimum Sufficiency Theorem}
(Theorem~\ref{thm:minimum-sufficiency}) establishes that the partition floor,
the minimum sufficient contact map, the minimum sufficient Lagrangian, and
the minimum sufficient action are in canonical bijective correspondence.
Every additional structural annotation is a partition depth inflation; none
reduces the floor.

Third, the \emph{Local Invisibility Theorem}
(Theorem~\ref{thm:local-invisibility}) shows that universal local validity
(every unit passes its local verifier) does not imply global contact map
coherence.  This is the formal basis for the claim that classical point-
evaluation testing cannot determine the contact map.

The \emph{Wind Tunnel DSL}, specified by a formal context-free grammar with
compositional, monotone, floor-aware operational semantics, provides the
language in which contact map investigations are specified, executed, and
recorded.  Its co-authorship model --- users write structural intuition, the
system writes investigation --- is a practical instantiation of the
theoretical framework.

The \emph{Water Level Invariance Theorem} (Theorem~\ref{thm:water-beads})
formalises the core epistemic insight: rearranging observable states (test
outcomes, type annotations, coverage reports) does not change the water level
$\bmin$.  The system's job is to measure $\bmin$ accurately, not to remove
it.

Future work includes: (i) constructive bounds on $\bmin$ for specific classes
of software systems (microservices, monoliths, functional pipelines); (ii) a
type-theoretic embedding of the DSL in a dependent type system; (iii)
empirical validation of the coordination regime boundaries against known
failure modes in production systems.

% ==========================================================================
% Bibliography
% ==========================================================================

\bibliographystyle{plainnat}
\bibliography{references}

\end{document}
